% Generated draft: requires mathematical and source review.
\documentclass[11pt]{article}
\usepackage[margin=25mm]{geometry}
\usepackage{fontspec}
\setmainfont{Latin Modern Roman}
\usepackage{amsmath,amssymb,mathtools,unicode-math}
\setmathfont{Latin Modern Math}
\usepackage{enumitem,needspace,xurl,hyperref}
\setlength{\parindent}{0pt}
\setlength{\parskip}{5pt}
\setlength{\emergencystretch}{3em}
\newcommand{\N}{\mathbb{N}}
\newcommand{\Z}{\mathbb{Z}}
\newcommand{\Q}{\mathbb{Q}}
\newcommand{\R}{\mathbb{R}}
\newcommand{\C}{\mathbb{C}}
\newcommand{\F}{\mathbb{F}}
\newcommand{\A}{\mathbb{A}}
\newcommand{\PP}{\mathbb{P}}
\newcommand{\E}{\mathbb{E}}
\newcommand{\eps}{\varepsilon}
\newcommand{\dd}{\,\mathrm{d}}
\begin{document}
\section*{UnsolvedMath ID: 30003615}
\textbf{Source problem code:} catalogue entry
\subsection{Definitions and mathematical statement}
The supplied record describes a research programme comprising three questions, rather than a single fully quantified conjecture. Its intended classes of tilings and subshifts, topology, and equivalence relations are not specified. The notation below makes these omissions explicit; the conventional models described here are conditional interpretations, not additional hypotheses attributed to the source.

For a tiling \(T\) of a Euclidean space \(\R^d\), where the dimension and permitted tiles remain unspecified, a patch is a collection of tiles of \(T\). Under a chosen adjacency rule, a corona about a distinguished tile consists of tiles reached within a prescribed number of adjacency steps. An aperiodic tiling, in the translational sense, satisfies
\[
\{v\in\R^d:T+v=T\}=\{0\}.
\]
The record does not specify whether this meaning of aperiodicity is intended, whether coronae must be finite, or which adjacency rule is used.

For topological spaces \(X\) and \(Y\), write
\[
\operatorname{Homeo}(X,Y)=\{h:X\to Y:h\text{ is bijective and both }h,h^{-1}\text{ are continuous}\}.
\]
A tiling space must be equipped with a topology before this notation can be applied; the construction and topology of the spaces in the question are not supplied.

One conventional model for a subshift is a closed, shift-invariant subset \(X\subseteq A^G\), with the product topology, where \(A\) is a finite discrete alphabet and \(G\) is a discrete group. For the convention \((\sigma^g x)(k)=x(g^{-1}k)\), one possible meaning of an extended symmetry is a pair
\[
(h,\alpha)\in\operatorname{Homeo}(X,X)\times\operatorname{Aut}(G)
\quad\text{such that}\quad
h\circ\sigma^g\circ h^{-1}=\sigma^{\alpha(g)}
\quad\text{for every }g\in G.
\]
The source does not specify the alphabet, acting group, or definition of extended symmetry, so this formula is only a possible interpretation.

\begin{enumerate}
\item Determine how aperiodic tilings can be reconstructed as limits of growing coronae. A precise formulation requires the class of tilings, the corona construction, the compatibility data retained between successive coronae, and the meaning of convergence. None is fixed by the supplied record.
\item Determine which homeomorphisms occur between the intended tiling spaces \(X\) and \(Y\), and how many there are. The source does not specify the spaces or whether the requested count concerns \(\operatorname{Homeo}(X,Y)\) itself or equivalence classes of such maps. It also does not impose compatibility with a dynamical action.
\item Determine which extended symmetries the intended subshifts can possess. The class of subshifts and the intended symmetry notion must be specified before this becomes a definite classification problem.
\end{enumerate}
These unresolved choices prevent recovery of a unique precise mathematical assertion from the supplied record.
\subsection{Short English statement}
The record asks how growing coronae reconstruct aperiodic tilings, which homeomorphisms occur between tiling spaces and how many there are, and which extended symmetries subshifts admit. It leaves the relevant classes, topology, convergence, counting convention, and symmetry definition unspecified. The displayed definitions give conditional interpretations; this statement-only draft requires human mathematical review.
\subsection{Sources}
\begin{itemize}
\item Supplied local UnsolvedMath registry record, permanent ID \(30003615\), title ``Topology and Symmetries of Aperiodic Tiling Spaces,'' exact source code \texttt{catalogue entry}, in the Oberwolfach Reports Open Problems collection. The project identifies \texttt{data/unsolvedmath/problems.json} and the supplementary \texttt{catalog/problems.json} as local registry resources; neither was independently inspected for this draft. The record gives proposed year 2017 and names no proposer.
\item Supplied report citation: \emph{Spectral Structures and Topological Methods in Mathematical Quasicrystals} (2018), \emph{Oberwolfach Reports}, pp.\ 2781--2845, DOI \texttt{10.4171/owr/2017/46}. Supplied source URL: \url{https://doi.org/10.4171/owr/2017/46}. The background attributes the questions to mathematical workshop discussions documented by Oberwolfach Reports. No individual contribution or problem-specific page is identified.
\item Supplied literature source: ``Oberwolfach Report 46 2017,'' \url{https://doi.org/10.4171/owr/2017/46}. The literature assessment checked at \texttt{2026-08-21} labels the questions open and states that no general resolution is known. This is supplied dated curation, not an independently verified literature assessment; the background explicitly requests independent verification of time-sensitive status and cited claims.
\item The registry labels the statement \texttt{corrected\_verified} and reports a manual semantic, typographic, and LaTeX review asserting self-containment and preservation of the mathematical target. These are supplied review claims, not verification performed here. The broad formulation still leaves the mathematical choices described above unresolved.
\end{itemize}
\end{document}
